\subsection{A native live-core instantiation}
\label{sec:live-core-instantiation}

The next-rekey comparison has a concrete record cost on the authoritative
live-session profile. Its application owns plaintext-bearing Send and
delivery; the network interface includes exact Resume, active control and
message processing, ACKs, status and public commitments. Storage follows
the native immutable-image and pending-intent history through failed writes,
process loss and reopening. Captures of old or unrelated records remain
available. A complete post-cut target-session acquisition invokes the
original eligibility stop. The selected successor-root and next-chain-seed
projections use the ownership and phase checks of the joint continuation.
The partial secret-field interface consists exactly of these two named,
phase-checked projections; other record acquisitions return the declared
complete objects. These observations and the private-application/Resume
boundary define the core's complete call surface. The retained live-profile
artifact records their source correspondence in
Appendix~\ref{app:joint-artifacts}.

This core obtains all its replies through the record, primitive and
application interfaces already constructed. Its ordinary calls expose
public metadata or declared available values. Consequently its charged
projection event is empty: $B_j=\varnothing$ for each candidate and
$\eta_{\rm proj}^{\rm nat}=0$. The report composition theorem supplies a
separate extension for cleanup and report interfaces with their declared
projection events.

For candidate $j$, let $M_j^{\rm seal}$ bound fresh outer sealing
invocations, including attempts whose subsequent commit fails. Let
$d_{j,i}$ bound the distinct private image/intent digest inputs in row $i$,
and let $Q_{j,i}$ bound the entire ordinary-$H$ computation of its derived
distinguisher. These are deterministic bounds fixed before execution and
cover every corresponding comparison world. Define
\begin{equation}
\begin{split}
 \varepsilon_{{\rm rec},j}
 &=\min\left\{1,
     \sum_{i=1}^{M_j^{\rm seal}}\Phi_{192}(d_{j,i},Q_{j,i})
       +\delta_{{\rm seal},j}\right\},\\
 \delta_{{\rm seal},j}
 &=\min\left\{1,\binom{M_j^{\rm seal}}{2}2^{-192}\right\}.
\end{split}
\label{eq:live-core-record-cost}
\end{equation}
The second term is the private-wire structural-equality comparison of
Corollary~\ref{cor:real-to-opaque}; it is zero when that equality interface
is absent. Copies and exact pending-target application retain their
original sealing handles. Atom-free rows have $d_{j,i}=0$.

\begin{corollary}[Next-rekey confidentiality on the live core]
\label{cor:live-core-next-rekey}
Use the native next-rekey and marked first-message interfaces of
Theorem~\ref{thm:cb-next-rekey} with the live core above, its recorded source
correspondence and finite whole-execution budgets. Then
\begin{equation}
\begin{split}
 \operatorname{Adv}_{\rm msg}\le\min\left\{\frac12,
 \frac12\left(\eta_{\rm attr}^{\rm nat}\right.\right.\\[-2pt]
 \left.\left.\qquad+
 \sum_{j\in J}\left(2\varepsilon_{{\rm rec},j}
      +a_j^++a_j^-+2\delta_j+B_j^{\rm cont}\right)
 \right)\right\}.
\end{split}
\label{eq:live-core-next-rekey}
\end{equation}
Here $B_j^{\rm cont}$ is the explicit ContextBound, root-initialization,
joint-schedule and AEAD score cost in~\eqref{eq:cb-continuation-bound}.
Attribution is evaluated in the original native experiment.
\end{corollary}
\begin{proof}
The live-profile representation invariant follows the original reservation,
saved-version and sealing handles through each native transition. On its
stopped interface, the materialized record machine has exactly the native
outputs and storage outcomes. Corollary~\ref{cor:real-to-opaque} and
Appendix~\ref{app:primitive-inputs} therefore give
$\rho_j\le2\varepsilon_{{\rm rec},j}$. The core's complete captures use
the existing acquisition stop, and its remaining replies use the declared
interfaces, so its charged projection event is empty. Substitute these two
facts into~\eqref{eq:cb-native-bound}.
\end{proof}

\paragraph{Source selection and attribution.}
The attribution term distinguishes two actual interfaces. In the general
experiment, the selector may name any original reservation handles with
the required session, epoch and role metadata. For example, a generation
reservation $g_1$ may be installed in a transaction-private
\texttt{Plan::Key} whose first persistence fails before a complete pending
target or image becomes durable. Reopening the prior
authoritative state permits a new reservation $g_2$, followed by an honest
accepted exchange using $(g_2,a_2)$. Choose
$c<\mathsf{res}(g_1)<\mathsf{res}(g_2)$ and
$c<\mathsf{res}(a_2)$, with the target signing keys and designated entropy
unexposed and with no acquisition of a post-cut target-session version.
A selector naming $(g_1,a_2)$ can pass
the reservation-event and recovery-cut checks while failing usable
attribution. Every control signature and transcript hash in this trace is
honest. This execution is already charged by
$\eta_{\rm attr}^{\rm nat}$ in the general theorem.

Section~\ref{sec:executed-source-attribution} defines an executed-slot
selector that registers $g,a$ before their first private KEM computations
return. A retained-plan induction preserves those original slots, and the
complete signed-body and wire bindings associate an accepted allocation
with them. Its signature experiment retains original signing randomizers
and exact resumptions. That selector receives the explicit signature/hash
bound in~\eqref{eq:executed-attribution-cost}; the general metadata selector
in~\eqref{eq:live-core-next-rekey} retains its original attribution term.
